% ---------------------------------------------------------------------------
% Chapter 5: The Children Get Married: Hadronisation and b-Fragmentation
% Generated from the outline; edit freely, this file is now the source.
% ---------------------------------------------------------------------------
\chapter{The Children Get Married: Hadronisation and $b$-Fragmentation}
\chaptermark{The Children Get Married}
\label{ch:children_getting_married}

\begin{journeybox}
No parton travels alone. Before the pilgrimage begins every child of the
family, and every cousin from the underlying event, has to marry: the colour
lines drawn by the shower and rearranged by colour reconnection become strings,
the strings break, and each quark leaves the house as one half of a meson or
one third of a baryon. The mother marries too, and the B~hadron she becomes
keeps most of her energy but not all of it.
\end{journeybox}

\physicsline{Lund string model, string breaking, $b$-fragmentation (Peterson, Lund--Bowler), $x_B$, hadron-level jet}

\section{From partons to hadrons: the Lund string}
\label{sec:children_getting_married:from_partons_to_hadrons_the_lund_string}
\todo{Write this section.}

\section{$b$-fragmentation: how much of the mother's energy the B~hadron keeps ($x_B$)}
\label{sec:children_getting_married:b_fragmentation_how_much_of_the_mother_s}
\todo{Write this section.}

\section{Our jet at this stage: the hadron-level family and the B~hadron's share}
\label{sec:children_getting_married:our_jet_at_this_stage_the_hadron_level_f}
\todo{Numbers, plots and event display of the $b$-jet at this stage.}

\begin{ledger}{The Children Get Married}
  \ledgerrow{before this stage}{}{}{}{--}
  \ledgerrow{after this stage}{}{}{}{}
\end{ledger}

\section{Further reading}
\label{sec:children_getting_married:further_reading}
Official and theory references for this stage: \cite{Andersson:1983ia}, \cite{Peterson:1982ak}, \cite{Bowler:1981sb}, \cite{ALEPH:2001pfo}.

\begin{pitfalls}
\todo{Pitfalls and FAQs for newcomers at this stage.}
\end{pitfalls}

\begin{exercises}
\item \todo{Exercise 1.}
\item \todo{Exercise 2.}
\end{exercises}
